% !TeX program = lualatex
% Compile twice: lualatex 125.tex
% All references and prose are embedded; no BibTeX database or external inputs.
\documentclass[11pt,a4paper]{article}
\usepackage[margin=24mm,headheight=14pt]{geometry}
\usepackage{fontspec}
\setmainfont{Latin Modern Roman}
\setsansfont{Latin Modern Sans}
\setmonofont{Latin Modern Mono}
\usepackage{amsmath,amssymb,mathtools}
\usepackage{unicode-math}
\setmathfont{Latin Modern Math}
\usepackage{microtype}
\usepackage{enumitem,xurl,needspace}
\usepackage[dvipsnames]{xcolor}
\usepackage{hyperref}
\hypersetup{unicode=true,colorlinks=true,linkcolor=MidnightBlue,urlcolor=MidnightBlue,
 pdftitle={500 Mathematical Problem Statements},pdfauthor={Source-annotated compilation},bookmarksnumbered=true}
\usepackage{bookmark}
\usepackage{tocloft}
\setlength{\cftsecnumwidth}{3em}
\cftsetpnumwidth{2.5em}
\cftsetrmarg{4em}
\usepackage{titlesec}
\titleformat{\section}{\Large\bfseries\raggedright}{\thesection}{0.7em}{}
\usepackage{fancyhdr}
\pagestyle{fancy}\fancyhf{}
\fancyhead[L]{\small\sffamily Mathematical problem statements}
\fancyhead[R]{\small Source edition 22}
\fancyfoot[C]{\thepage}
\setlength{\parindent}{0pt}
\setlength{\parskip}{5pt plus 1pt}
\setlength{\emergencystretch}{3em}
\setcounter{tocdepth}{1}
\setcounter{secnumdepth}{1}
\setlist[itemize]{leftmargin=3em,itemsep=5pt,topsep=4pt,parsep=0pt}
\allowdisplaybreaks
\newcommand{\N}{\mathbb{N}}
\newcommand{\Z}{\mathbb{Z}}
\newcommand{\Q}{\mathbb{Q}}
\newcommand{\R}{\mathbb{R}}
\newcommand{\C}{\mathbb{C}}
\newcommand{\F}{\mathbb{F}}
\newcommand{\A}{\mathbb{A}}
\newcommand{\PP}{\mathbb P}
\newcommand{\E}{\mathbb{E}}
\newcommand{\eps}{\varepsilon}
\newcommand{\dd}{\,\mathrm d}
\newcommand{\sref}[2]{\hyperlink{source:#1:#2}{[S#2]}}
\newcommand{\eref}[2]{\hyperlink{extra:#1:#2}{[E#2]}}
% Turn the next switch off for a shorter reading copy. The quotations preserve
% every exactTarget field, including qualifications omitted from a short summary.
\newif\ifshowcatalogue
\showcataloguetrue
\newcommand{\cataloguescope}[1]{\ifshowcatalogue
 \paragraph{Exact scope in the supplied catalogue.}
 \begin{quote}\small #1\end{quote}\fi}
\hypersetup{pdftitle={Problem 125 - Polynomial bound for limit cycles},pdfauthor={Open Problems Math research notebook}}
\fancyhead[L]{\small\sffamily Open Problems Math \textbar\ Research notebook}
\fancyhead[R]{\small\sffamily 125 \textbar\ 27 September 2026}
\numberwithin{equation}{section}
\begin{document}
\setcounter{section}{124}
\setlength{\parskip}{3pt plus 1pt}
\setlist[itemize]{itemsep=3pt,topsep=3pt}
\titlespacing*{\paragraph}{0pt}{1.5ex plus .5ex minus .2ex}{1em}
% ================================================================
% problemId: problem.polynomial-bound-for-limit-cycles-of-planar-polynomial-vector-fields
% source releaseRank: 125; source releaseStatus: open
% exactTarget SHA-256: 6925cd0b4a6199e6ab0c77eea04000a2af23ea09981a423debaedf3425dd91b3
\clearpage
\section{\texorpdfstring{Polynomial Bound for Limit Cycles of Planar Polynomial Vector Fields}{Polynomial Bound for Limit Cycles of Planar Polynomial Vector Fields}}\label{125}
\subsection{Definitions and mathematical statement}
For real polynomials $P,Q$ of degree at most $d$, a limit cycle of $\dot x=P(x,y)$, $\dot y=Q(x,y)$ is an isolated periodic orbit. The target asks whether there is one exponent $q>0$ with
\[
 \#\{\text{limit cycles of }(P,Q)\}\le d^q
\]
for every degree $d$ and every such system, with the nontrivial degree range $d\ge2$. The exponent is independent of all coefficients and of $d$. A finite but faster-than-polynomial degree bound would not answer this stronger quantitative question.
\par\smallskip\textit{Formulation and concept references:} \sref{125}{1}.
\cataloguescope{Determine whether the number K of limit cycles of every planar polynomial differential system of degree at most d satisfies K \ensuremath{\le} d\textasciicircum{}q for some universal exponent q.}
\subsection{Short English statement}
Can the number of limit cycles of every planar polynomial system be bounded by one universal power of its degree?
% BEGIN RESEARCH ADDITION 125
\subsection{Research attempt}

\paragraph{Current frontier.}
Smale's Problem 13 asks for the universal polynomial degree bound in the
catalogue, not merely finiteness for a fixed vector field.
\sref{125}{1}, Problem 13.
Buzzi--Novaes' 2024 analysis documents why a proposed quadratic formula
for the Hilbert numbers conflicts with established lower constructions
of order $d^2\log d$.  That refutation does not refute a bound $d^q$
with some larger universal exponent. \eref{125}{1}, Sections 1--3.
The present attempt actively constructs systems with many isolated
cycles, keeps track of their degree, and checks whether the construction
can outgrow every polynomial.

\paragraph{Attack routes.}
A counterexample could come from many separated oscillating regions,
from many nested cycles, or from a recursive construction whose cycle
count grows faster than the cost in degree.  A proof would need uniform
control of zeros of return maps as coefficients vary and degenerate.
Counting zeros of a polynomial surrogate is insufficient unless an exact
relation with isolated periodic orbits is established.  We combine the
first two construction mechanisms explicitly, then calculate why that
combination alone fails to give superpolynomial growth.

\paragraph{Attempt I: many wells with controlled algebraic degree.}
Let $r\geq1$, and choose a real polynomial $f$ of degree $r$ with $r$
distinct real roots, for example $f(t)=\prod_{i=1}^r(t-i)$.  Put
\[
 H(x,y)=f(x)^2+f(y)^2.
\]
There is $a>0$ such that $f^{-1}((-a,a))$ consists of $r$ disjoint
intervals, each containing one root, and $f$ maps each interval
smoothly and bijectively onto $(-a,a)$.  To justify this globally,
choose disjoint inverse-function neighborhoods of the simple roots.
The proper function $|f|$ has a positive minimum on their complement
after slightly shrinking their images; choose $a$ below that minimum.
Thus no additional small-value component is omitted.

On each of the $r^2$ resulting rectangles, the map
$(x,y)\mapsto(f(x),f(y))$ is a diffeomorphism onto $(-a,a)^2$.
For every $0<c<a^2$, the level $H=c$ is therefore exactly $r^2$
disjoint smooth closed ovals, one in each rectangle.  There are no other
components, since $H<c$ forces both $|f(x)|,|f(y)|<a$.  Its gradient
is nonzero on every such level: in these local coordinates the level
is a nonzero circle and the coordinate change is invertible.

Choose distinct values $0<c_1<\cdots<c_m<a^2$ and write
$h(s)=\prod_{j=1}^m(s-c_j)$.  For the quarter-turn matrix
$J=\begin{pmatrix}0&-1\\1&0\end{pmatrix}$ define the polynomial vector
field
\begin{equation}\label{125:field}
 V=J\nabla H-h(H)\nabla H.
\end{equation}
Its degree is exactly
\begin{equation}\label{125:degree}
 d=2r(m+1)-1.
\end{equation}
The highest-degree part comes from $h(H)\nabla H$; it is nonzero and
has strictly greater degree than $J\nabla H$.
Most importantly,
\begin{equation}\label{125:energy}
 \dot H=\nabla H\cdot V=-h(H)|\nabla H|^2.
\end{equation}
Every oval in $H=c_j$ is invariant.  On it the vector field equals
$J\nabla H$, which is everywhere nonzero.  A smooth nonvanishing
vector field on a compact circle traverses it in finite period,
so these are periodic orbits, not merely invariant sets.

\paragraph{Attempt II: prove isolation and count all cycles.}
Use the transverse coordinate $s=H-c_j$ near one oval.  Linearizing
\eqref{125:energy} along its period-$T_j$ trajectory $\gamma$ gives
\[
 \dot s=-h'(c_j)|\nabla H(\gamma(t))|^2s.
\]
Its return multiplier is
\begin{equation}\label{125:multiplier}
 \exp\!\left(-h'(c_j)\int_0^{T_j}
                      |\nabla H(\gamma(t))|^2dt\right)\ne1.
\end{equation}
The integral is strictly positive and $h'(c_j)\ne0$.  Thus every
constructed periodic orbit is hyperbolic and isolated, and is genuinely
a limit cycle in the definition of the problem.

There are no further nonconstant periodic orbits.  A nonstationary
trajectory cannot hit a critical point of $H$, because there $V=0$
and uniqueness of the smooth ordinary differential equation would make
it stationary.  Between consecutive roots of $h$, \eqref{125:energy}
then has a fixed strict sign.  A trajectory cannot cross a level
$H=c_j$ either: that regular level is invariant and uniqueness prevents
crossing.  Consequently a periodic trajectory must have constant $H$
equal to one of the $c_j$, and every such level has already been counted.
We have proved an exact count
\begin{equation}\label{125:count}
 K=mr^2
\end{equation}
for this family.  This conclusion is analytic, not inferred from a phase
portrait or a numerical orbit integrator.

For a concrete example take $f(t)=t(t-1)$ and $c_1=1/100$.
The map has four wells, and the above construction gives four hyperbolic
limit cycles for a degree-seven field.  The threshold can be checked
explicitly: the only nonzero critical value of $f$ in magnitude is
$1/4$, so one may take any $a<1/4$ with $c_1<a^2$.
The companion script expands this example and verifies
\eqref{125:energy} and the degree exactly.

\paragraph{Attempt III: can nesting beat the degree cost?}
Combining \eqref{125:degree} and \eqref{125:count} gives
\[
 K=\frac{m}{4(m+1)^2}(d+1)^2
       \leq\frac{(d+1)^2}{16},
\]
since $(m+1)^2\geq4m$.  Increasing the number of nested levels at each
well consumes degree too quickly; this entire construction remains
quadratic in degree.  The calculation is a barrier for this chosen
ansatz, not a universal upper bound.  In particular, it cannot compete
with the known $d^2\log d$ lower constructions recorded in
\eref{125}{1}.

One may try to insert oscillatory factors with many real roots but low
polynomial degree.  For a genuine polynomial $h$, its number of distinct
roots is at most its degree, which is exactly the cost in
\eqref{125:degree}.  Replacing it by an oscillatory transcendental
function would leave the problem's polynomial class.  Alternatively,
using many algebraic wells while pretending that their defining
polynomial has fixed degree contradicts the same degree bookkeeping.
These explicit failures explain why this counterexample search stops
without suggesting that every possible construction is similarly bounded.

On the proof side, the exact identity \eqref{125:energy} is not available
for a general polynomial vector field.  Its Poincar\'e return map is
analytic on its domain, but need not be a polynomial of degree controlled
by $d$.  Analytic functions can have arbitrarily many isolated zeros as
their coefficients vary.  To apply a zero-counting theorem one would
need a uniform complexity bound for the actual return maps, including
near separatrices and singular limiting systems.  No such general bound
is established here.

\paragraph{Stress test.}
The construction counts isolated nonconstant periodic orbits, not
centers, equilibria, or arbitrary compact algebraic components.  Its
multiplier verifies robustness under sufficiently small perturbations
of a fixed constructed field, but the allowable perturbation size may
depend on $r,m$ and the selected levels.  That local robustness does not
produce a uniform theorem over all coefficient space.

The catalogue permits any universal exponent $q>0$ and no coefficient-
height dependence.  A lower bound exceeding every quadratic polynomial
rules out $q\leq2$, not all larger $q$.  Conversely, per-field finiteness
and normalization of the coefficient vector do not by themselves give
a uniform bound: the normalized parameter space contains degenerate
limits, and the cycle-counting function need not be bounded by a
naive compactness argument.  The critique in \eref{125}{1} illustrates
why changing the definition of a cycle to a computable surrogate would
invalidate the target rather than solve it.

\paragraph{Outcome.}
\textbf{Outcome: Exploratory approach with unresolved gap.}

A two-parameter polynomial family with exactly $mr^2$ hyperbolic limit
cycles and degree $2r(m+1)-1$ has been constructed and proved to have
no other periodic cycles.  Its degree accounting shows that separated
wells plus nested algebraic levels cannot yield a superpolynomial
counterexample by this mechanism.  The general return-map bound or a
more efficient recursive counterexample construction remains missing.
This family is not claimed as a new lower-bound record or a resolution.
% END RESEARCH ADDITION 125
\subsection{Sources}
\begin{itemize}\raggedright
\item[\textnormal{[S1]}]\hypertarget{source:125:1}{} Steve Smale, Mathematical Problems for the Next Century, 1998.
\url{https://www.fim.uni-passau.de/fileadmin/dokumente/fakultaeten/fim/lehrstuhl/muller/SmaleProblems1998.pdf}.
{\small\textit{Catalogue formulation source.}}
% BEGIN RESEARCH REFERENCES 125
\item[\textnormal{[E1]}]\hypertarget{extra:125:1}{} Claudio A. Buzzi and Douglas D. Novaes.
\emph{A note on a recent attempt to solve the second part of Hilbert\textquotesingle s 16th Problem}, arXiv:2411.09594, 14 November 2024. Discussion of the claimed quadratic formula and the established superquadratic lower constructions.
\url{https://arxiv.org/html/2411.09594}.
{\small\textit{Consulted 27 September 2026 for the indicated result or scope.}}
% END RESEARCH REFERENCES 125
\end{itemize}


\end{document}
